\chap{Abstract}

This dissertation studies how heterogeneity in financial environments shapes
household wealth accumulation and realized investment performance. The first
essay develops a heterogeneous-agent model with dispersion in household asset
returns and shows that return heterogeneity substantially improves the model's
ability to match empirical wealth moments. It also compares the distributional
and welfare implications of wealth taxation and capital income taxation when
households earn different returns. The second essay uses the RAND
Longitudinal HRS to measure returns to net wealth and documents a hump-shaped
relationship between generalized trust and realized investment performance.
Across cross-sectional and panel specifications, moderate trust is associated
with the highest returns even after accounting for demographic characteristics
and portfolio composition. Taken together, the essays show that persistent
differences in returns, beliefs, and financial environments are central to
understanding wealth inequality and household financial behavior.

\keywords{\DissertationKeywords}

\begin{singlespace}
\section*{Primary Reader and Dissertation Advisor}
\PrimaryReader \\
Department of Economics \\
Johns Hopkins University

\section*{Secondary Readers}
\SecondaryReaderOne \\
Department of Economics \\
Johns Hopkins University

\vspace{0.1in}

\SecondaryReaderTwo \\
Department of Economics \\
Johns Hopkins University
\end{singlespace}
